\chapter{Referencial Te\'orico}
\label{cap:referencial}

Texto introdut\'orio de um par\'agrafo, conforme exigido pela ABNT. Esta
parte do trabalho demanda o dom\'inio das regras de cita\c{c}\~ao (diretas
e indiretas). Reveja as regras estabelecidas pela ABNT (NBR 10520) para
cita\c{c}\~oes diretas e indiretas e indica\c{c}\~ao de fontes de gr\'aficos,
figuras ou outro recurso ilustrativo utilizado.

\section{T\'itulo representativo do primeiro conceito}
\label{sec:conceito1}

Discuss\~ao espec\'ifica sobre a articula\c{c}\~ao das informa\c{c}\~oes com
foco na delimita\c{c}\~ao proposta.

\subsection{T\'itulo representativo de um subt\'opico}

Exemplo de cita\c{c}\~ao direta de at\'e 3 linhas, que permanece no corpo do
texto, entre aspas: de acordo com \textcite{sa1971}, isso acontece ``por meio
da mesma arte de conversa\c{c}\~ao que abrange t\~ao extensa e significativa
parte da nossa exist\^encia cotidiana''.

Exemplo de cita\c{c}\~ao direta com mais de 3 linhas (ambiente
\texttt{citacao}, recuo de 4\,cm, fonte tamanho 10, espa\c{c}amento simples,
sem aspas), conforme \textcite{nicholls1993}:

\begin{citacao}
A teleconfer\^encia permite ao indiv\'iduo participar de um encontro
nacional ou regional sem a necessidade de deixar seu local de origem.
Tipos comuns de teleconfer\^encia incluem o uso da televis\~ao, telefone e
computador. Atrav\'es de \'audio-confer\^encia, utilizando a companhia
local de telefone, um sinal de \'audio pode ser emitido em um sal\~ao de
qualquer dimens\~ao. \parencite[p.~181]{nicholls1993}
\end{citacao}

Exemplo de cita\c{c}\~ao indireta: para \textcite{authier1982}, a ironia
seria assim uma forma impl\'icita de heterogeneidade.

Exemplo de inser\c{c}\~ao de figura com fonte obrigat\'oria:

\begin{figure}[htbp]
  \centering
  % \includegraphics[width=0.7\textwidth]{figuras/exemplo.png}
  \fbox{\parbox{0.7\textwidth}{\centering\vspace{2cm}Espa\c{c}o reservado para a figura\vspace{2cm}}}
  \caption{T\'itulo da figura}
  \label{fig:exemplo}
  \fonte{Do autor (2026).}
\end{figure}

Para mais informa\c{c}\~oes sobre cita\c{c}\~ao, consulte a norma NBR 10520:2002.
